
%(BEGIN_QUESTION)
% Copyright 2016, Tony R. Kuphaldt, released under the Creative Commons Attribution License (v 1.0)
% This means you may do almost anything with this work of mine, so long as you give me proper credit

En relativt ny teknologi for å måle konsentrasjonen av oksygengass oppløst i vann kalles {\it dynamic luminescence quenching}. Dette er en optisk teknologi som bygger på fluorescensegenskapene til et stoff som inneholder "lumifor"-molekyler. Grunntanken er: lumifor-molekylene fluorescerer lett når de eksponeres for en bestemt lysbølgelengde, men fluorescensen påvirkes av tilstedeværelsen av oksygenmolekyler. Oksygen nær lumifor-molekylene har en tendens til å "slokke" den eksiterte tilstanden i disse molekylene etter eksponering for innfallende lys, og hindrer dermed fluorescens. Ved å eksitere lumifor-molekylene gjentatte ganger med lyspulser og måle det fluorescerende lyset som returneres fra molekylene, kan konsentrasjonen av oppløst oksygen estimeres.

\vskip 10pt

Ulike lysfarger brukes i denne analysertypen: én farge sendes inn i lumifor-molekylene fra en lyskilde, og en annen farge sendes ut fra lumifor-molekylene og detekteres av en lyssensor. Den ene av disse fargene er {\it blå} og den andre er {\it rød}.

\vskip 10pt

Basert på det du vet om samspillet mellom lys og molekyler, identifiser hvilken farge som sendes ut av lyskilden og hvilken som sendes ut av lumifor-molekylene. Bestem også om intensiteten i lyset som sendes ut av lumifor-molekylene er {\it direkte proporsjonal} med eller {\it omvendt proporsjonal} med konsentrasjonen av oppløst oksygen i væskeprøven.

\underbar{file i00635no}
%(END_QUESTION)





%(BEGIN_ANSWER)

Lyskilden sender ut blått lys, mens lumifor-molekylene fluorescerer rødt lys. Det vet vi fordi innkommende bølgelengde alltid er kortere (høyere frekvens) enn bølgelengden ved fluorescens. Blått lys har kortere bølgelengde (høyere frekvens) enn rødt lys, og derfor må blått være fargen fra lyskilden og rødt være fargen som returneres fra lumifor-molekylene når de fluorescerer.

\vskip 10pt

Proporsjonaliteten mellom intensiteten i dette fluorescerende lyset og konsentrasjonen av oppløst oksygen er {\it omvendt}, basert på beskrivelsen av dynamic luminescence quenching i starten av oppgaven. Jo mer oppløst oksygen i nærheten av lumifor-molekylene, desto mer slokkes fluorescensen deres, og desto mindre rødt lys vil de sende ut når de eksiteres av blått lys fra lyskilden.

%(END_ANSWER)





%(BEGIN_NOTES)


%INDEX% Måling, analytisk: oppløst oksygen
%INDEX% Fysikk, optikk: dynamic luminescence quenching
%INDEX% Fysikk, optikk: fluorescens

%(END_NOTES)
